%!TEX root=../main.tex
\section{Feasibility}
\label{sec:feasibility}

We now show Calibration-Preserving Parity is a feasible definition
of fairness through a proof by construction. We will show that it is possible
derive calibrated classifiers with minimal disparity from classifiers that do
not meet this condition.

Assume we are given classifiers $h_1$ and $h_2$ for groups $G_1$ and $G_2$.
$h_1$ and $h_2$ do not achieve the disparity condition, but are approximately
calibrated --- i.e.  both $\calib{h_1}$ and $\calib{h_2}$ are small.\footnote{
  If this is not the case, one could improve calibration
  with Platt Scaling \cite{platt1999probabilistic} or
  isotonic regression \cite{zadrozny2001obtaining}.
}
Without loss of generality, for the remainder of the paper we assume that
\fairconst{h_1}{1} is larger than \fairconst{h_2}{2}.
We will show that, under minor assumptions, that is is possible to achieve
Calibration-Preserving Parity with respect to $\fairconstname$:
%
\begin{thm}[Feasibility]
  Let $h_1$ and $h_2$ be defined as above. If $\fairconst{h_1}{1} \leq
  \argmax_{h^*_2 \in \mathcal{H}^*_2} \fairconst{h^*_2}{2}$, where
  $\mathcal{H}^*_2$ is the set of calibrated classifiers for group $G_2$,
  then there exists some classifier $\fair h_2$ for group $G_2$ with
  $\calib{\fair h_2} \leq \calib{h_2}$ and $\fairconst{h_1}{1} = \fairconst{\fair h_2}{2}$.
  \label{thm:feasibility}
\end{thm}

To show that this is possible, we will derive a new classifier $\fair h_2$ that
replaces some predictions of $h_2$ with predictions from a ``trivial
classifier'' $h^\triv_2$ that simply outputs the base rate $\br_2$
on all inputs (i.e. $\forall \x$, $h^\triv_2(\x) = \br_2$).
Specifically, for some Bernoulli variable $B$ with expectation $\alpha$,
%
\begin{equation}
  \fair{h}_2(\x) = (1 - B) h_2(\x) + B h^\triv_2(\x).
  \label{eqn:fair_classifier}
\end{equation}
%
Intuitively, by modifying the parameter $\alpha$, we are interpolating between the two classifiers
$h_2$ and $h^\triv_2$, and thus interpolating between their costs.
A geometric interpretation of this can be seen in \autoref{fig:postprocess}.
By modifying the replacement rate $\alpha$, we can raise and lower the cost of $\fair h_2$.
Setting
$ \alpha = \frac{\fairconst{h_1}{1} - \fairconst{h_2}{2}}{\fairconst{h^\triv_2}{2} - \fairconst{h_2}{2}}$
ensures that the cost of $\fair h_2$ matches that of $h_1$ (see \autoref{thm:interpolation}
in \autoref{sec:detailed-feasibility}).
Moreover, the resulting $\fair h_2$ will be at least as calibrated as $h_2$
(see \autoref{thm:calibration} in \autoref{sec:detailed-feasibility}).
Therefore, $h_1$ and $\fair h_2$ will satisfied
Calibration-Preserving Parity with respect to \fairconstname.
We summarize these steps in \autoref{alg:method}, and prove that it meets the
conditions of \autoref{def:fair} in the supplementary material.
%
\begin{algorithm}
  \caption{\methodname{} by post-processing}
  \label{alg:method}
  \begin{algorithmic}
    \STATE {\bfseries Input:} classifiers $h_1$ and $h_2$ s.t.
    $\fairconst{h_2}{2} \leq \fairconst{h_1}{1} \leq \fairconst{h^\triv_2}{2}$, holdout set $P_{valid}$

    \begin{itemize}
      \STATE Determine base rate $\br_2$ of $G_2$ (using
        $P_{valid}$) to produce trivial classifier $h^\triv_2$

      \STATE Construct $\fair h_2$ using \eqref{eqn:fair_classifier} with
        $ \alpha = \frac{\fairconst{h_1}{1} - \fairconst{h_2}{2}}{\fairconst{h^\triv_2}{2} - \fairconst{h_2}{2}}$,
        where $\alpha$ is the interpolation parameter.
    \end{itemize}

    \RETURN $h_1$, $\fair h_2$.
  \end{algorithmic}
\end{algorithm}

\paragraph{Limitations.}
%
It is worth noting the cases in which \methodname{} is infeasible. For the following cases,
we assume that we cannot improve the cost $h_1$ and $h_2$ with further training
or postprocessing, given a finite set
of data and features.\footnote{
  Otherwise, the improved classifier would then be the input to \autoref{alg:method}.
}
%
\begin{corollary}
  \methodname{} is infeasible if
  $\fairconst{h_1}{1} > \argmax_{h^*_2 \in \mathcal{H}^*_2} \fairconst{h^*_2}{2}$.
  %where $\mathcal{H}^*_2$ is the set of all calibrated classifiers for group $G_2$.
  \label{crly:infeasible_trivial}
\end{corollary}
%
\begin{corollary}
  \methodname{} is infeasible if group membership $G_t$ is more predictive
  than any features --- i.e., if
  $ \fairconst{h^\triv_t}{t} \leq \fairconst{h'_t}{t}$, where $h'_t$ is
  any (possibly uncalibrated) classifier.
  \label{crly:infeasible_predictive}
\end{corollary}
%
\begin{corollary}
  In all other cases, \autoref{alg:method} will achieve \methodname{}.
  \label{crly:method_works}
\end{corollary}
%
(Recall that $\mathcal{H}^*_t$ is the set of calibrated classifiers for group $G_t$.)
\autoref{crly:infeasible_trivial} is trivially true because we cannot find a
calibrated classifier for group $G_2$ with cost equal to $\fairconst{h_1}{1}$.
\autoref{crly:infeasible_predictive} is a direct consequence of this result, combined with the fact that
the trivial classifier $h^\triv_t$ is the calibrated classifier for group $G_t$ that
maximizes any cost function of the form \eqref{eqn:cost_general}
(see \autoref{lma:trivial} in \autoref{sec:detailed-feasibility}).
To see why \autoref{crly:method_works} holds, note that \autoref{alg:method}
can match the cost of any classifier in the range $\left[ \fairconst{h_2}{2}, \fairconst{h^\triv_2}{2} \right]$
(see \autoref{thm:interpolation}).
Because $h^\triv_2$ maximizes cost, \autoref{alg:method} will fail only if
$\fairconst{h_1}{1} \geq \fairconst{h^\triv_2}{2} = \argmax_{h^*_2 \in
\mathcal{H}^*_2} \fairconst{h^*_2}{2}$, in which case \methodname{} is
infeasible.

%\paragraph{Relation to Equalized Odds via post-processing}

%\autoref{alg:method} provides an interesting connection to the work of
%\citet{hardt2016equality}, who
%provide a post-processing method to achieve Equalized Odds (false-positive
%and false-negative constraints without calibration).
%Their method amounts to a linear program over two dimensions, which is exactly
%what one would expect given that they are optimizing over a four-dimensional
%space with two linear constraints, leaving two free variables to optimize over.
%\autoref{alg:method} is a linear program over one dimension, as calibration
%constraints the set of classifiers to a one-dimensional space.
%\gp{This maybe goes to supplementary\ldots}


